\section{Design Decisions and Next Steps}

\subsection{Vector Embedding Decision}

Vector embedding is \textbf{not required} for the authoritative drug-drug interaction check. The safety path should use reviewed ingredient identity, accepted component-to-safety mappings, canonical interaction pairs, and relational joins. A similarity score cannot establish that two drug identities are equivalent and must not independently generate or suppress a medical warning.

Embeddings may be added later as an auxiliary retrieval feature for:
\begin{itemize}
  \item proposing medication or ingredient candidates from noisy OCR text;
  \item retrieving similar unresolved names for reviewer triage; or
  \item searching evidence passages semantically.
\end{itemize}

If introduced, embeddings should be stored separately with model name, model version, source-text hash, creation time, and entity reference. Candidate retrieval must still pass deterministic validation and, for safety mappings, medical review. This preserves a clean boundary between probabilistic assistance and operative clinical data.

\subsection{Recommended Physical Platform}

PostgreSQL is a suitable implementation target because the design benefits from UUIDs, robust constraints, transactional integrity, JSON support for preserved source payloads, partial indexes, and optional row-level security. This is a recommendation rather than a committed project constraint. The physical design should standardize UUID generation, timestamp/timezone handling, controlled status values, naming conventions, and migration tooling before schema creation.

\subsection{Open Design Decisions}

\begin{longtable}{L{5.0cm}L{6.2cm}L{4.0cm}}
\caption{Open decisions before physical schema approval}\label{tab:open-decisions}\\
\toprule
\textbf{Decision} & \textbf{Options/impact} & \textbf{Required owner}\\
\midrule
\endfirsthead
\toprule
\textbf{Decision} & \textbf{Options/impact} & \textbf{Required owner}\\
\midrule
\endhead
Database engine and deployment topology & PostgreSQL version, managed/self-hosted, environments, HA, backup objectives & Technical lead\\
Patient identity and minimum profile & Required demographics, identity verification, consent, account recovery & Product owner and privacy reviewer\\
Professional/reviewer qualification & License fields, verification process, review scopes, second-review thresholds & Clinical owner and administrator\\
Dataset release policy & Approval quorum, rollback, supersession, active channels, effective dates & Medical data owner\\
Interaction severity vocabulary & Source-specific levels versus normalized severity and conflict handling & Clinical reviewer\\
Retention and deletion & Intake documents, messages, tickets, audit, evidence, and health records & Privacy/legal owner\\
Notification provider and delivery policy & Channels, retry, quiet hours, timezone, escalation, idempotency & Product and technical leads\\
Physical column dictionary & Types, nullability, defaults, checks, unique constraints, and indexes for all 64 tables & Database designer\\
\bottomrule
\end{longtable}

\subsection{Implementation Sequence}

\begin{enumerate}
  \item Approve entity boundaries, role responsibilities, review workflow, and release semantics in this logical model.
  \item Produce a physical data dictionary for all tables, including types, nullability, constraints, and state values.
  \item Implement access control, dataset governance, and reference catalog first because later modules depend on them.
  \item Build a repeatable CSV staging/import pipeline with validation reports and reviewer work queues.
  \item Load a non-operative draft dataset and verify row counts, unresolved records, hashes, and source traceability.
  \item Add the separate interaction-pair dataset and evidence before enabling safety assessment.
  \item Implement patient workflows incrementally: intake, identification, snapshot/assessment, HCP review, then regimen/reminder.
  \item Run integrity, authorization, migration, backup/restore, performance, and clinical-review acceptance tests before any production release.
\end{enumerate}

\subsection{Conclusion}

The proposed model separates patient operations from governed reference data while preserving the joins needed for end-to-end medication safety. Its most important control is the boundary between candidate data and released, operative data: OCR output, indication text, and DDInter mappings remain reviewable evidence until accepted by an authorized reviewer. With the remaining physical decisions resolved, the 64-entity logical model can be converted into a normalized SQL schema without changing the core workflows documented in PA2.

